\section{Discussion and limitations}
\label{sec:discussion}
The findings connect the diversity of the surrogate library to a practical
model choice procedure. Different representations and uses of coarse data
provide accurate predictions on different problems. Reserved fine examples
then guide how much each model contributes before the next field is known.
Mixtures often improve on selecting one model, while inverse error weighting
shows that much of this gain is accessible through individual error estimates.
For an application with expensive simulations, additional surrogate computation
can therefore be a useful way to exploit the available fields. Once the library
is trained, comparing and fitting the three rules reuses its predictions without
retraining the models.

These results do not isolate a benefit specific to multifidelity data. The library includes
a fine only model, many datasets have equally sized coarse and fine training
tables, and validation allocations vary. Establishing data efficiency requires a
fine only library and alternative uses of fitting labels under equal total
label and computation budgets, counting every candidate and shared coarse
ensemble. Whether the additional surrogate computation is worthwhile depends
on simulation cost, the number of prediction queries, and the required accuracy.
Net cost savings and comparison with complete engineering automation systems
remain open.

\section{Conclusion}
The strongest surrogate for multifidelity field prediction changes across
problems. \method{} combines an adapted library of nine neural and reduced
basis models with an ensemble fitting workflow. Reserved fine solutions
determine which models to select or combine, and new inputs require no
additional coarse simulation. The assembled benchmark compares these choices
across field structures and fidelity relationships.

Across 22 datasets, the fitted mixture improves on selecting one model on 16
datasets using five reserved fine examples. Inverse error weighting captures most of this benefit.
ERA5 extends the comparison beyond PDE simulations. These findings support ensemble
fitting when uncertain model choice and expensive field data justify additional
surrogate computation.
